% Part: computability
% Chapter: recursive-functions
% Section: pr-functions-computable

\documentclass[../../../include/open-logic-section]{subfiles}

\begin{document}

\olfileid{cmp}{rec}{cmp}
\olsection{आदिम पुनरावर्ती फलने संगणनक्षम असतात}

फलन $h$ हे आदिम पुनरावर्तनाने परिभाषित केले आहे असे समजा:
\begin{eqnarray*}
h(\vec x, 0)     & = & f(\vec x) \\
h(\vec x, y + 1) & = & g(\vec x, y, h(\vec x, y))
\end{eqnarray*}
आणि $f$ व $g$ ही फलने संगणनक्षम आहेत असे समजा. ($x_0$, \dots,
$x_{k-1}$ यांचे संक्षिप्त रूप म्हणून आपण $\vec x$ वापरतो.) मग
$h(\vec x, 0)$ संगणित करता येते हे उघड आहे, कारण ते फक्त $f(\vec x)$
आहे आणि ते संगणनक्षम आहे असे आपण गृहीत धरले आहे. त्यानंतर
$h(\vec x, 1)$ देखील संगणित करता येते, कारण $1 = 0 + 1$ आणि म्हणून
$h(\vec x, 1)$ हे फक्त पुढील फलन आहे:
\begin{align*}
 h(\vec x, 1) & = g(\vec x, 0, h(\vec x, 0)) =  g(\vec x, 0, f(\vec x)).
\intertext{याच प्रकारे पुढे जाऊन आपण पुढील मूल्ये संगणित करू शकतो:}
h(\vec x, 2) & = g(\vec x, 1, h(\vec x, 1)) = g(\vec x, 1, g(\vec x, 0, f(\vec x)))\\
h(\vec x, 3) & = g(\vec x, 2, h(\vec x, 2)) = g(\vec x, 2, g(\vec x, 1, g(\vec x, 0, f(\vec x))))\\
h(\vec x, 4) & = g(\vec x, 3, h(\vec x, 3)) = g(\vec x, 3, g(\vec x, 2, g(\vec x, 1, g(\vec x, 0, f(\vec x)))))\\
& \vdots
\end{align*}
अशा रीतीने, सर्वसाधारणपणे $h(\vec x, y)$ संगणित करण्यासाठी
$h(\vec x, 0)$, $h(\vec x, 1)$, \dots, ही मूल्ये क्रमाक्रमाने
$h(\vec x, y)$ पर्यंत संगणित करा.

म्हणून $f$ आणि $g$ ही फलने संगणनक्षम असतील, तर आदिम पुनरावर्ती
व्याख्येपासून नवीन संगणनक्षम फलन मिळते. त्याचप्रमाणे $f$ आणि $g_i$ ही
फलने संगणनक्षम असतील, तर त्यांच्या संयोजनापासूनही संगणनक्षम फलन मिळते.

मूलभूत फलने $\Zero$, $\Succ$ आणि $\Proj{n}{i}$ ही संगणनक्षम आहेत, आणि
संयोजन व आदिम पुनरावर्तन यांमुळे संगणनक्षम फलनांपासून संगणनक्षम फलने
मिळतात. म्हणून प्रत्येक आदिम पुनरावर्ती फलन संगणनक्षम असते.

\end{document}
